\documentclass[11pt,a4paper]{article}

% ------------------------------------------------------------------
% Bayesian Learning and Monte Carlo Simulations -- Final Project 2026
% Project 11: Hourly City Load Factor (Tetouan)
% ------------------------------------------------------------------

\usepackage[utf8]{inputenc}
\usepackage[T1]{fontenc}
\usepackage[margin=2.5cm]{geometry}
\usepackage{amsmath,amssymb,amsfonts}
\usepackage{bm}
\usepackage{graphicx}
\usepackage{booktabs}
\usepackage{array}
\usepackage{caption}
\usepackage{subcaption}
\usepackage{float}
\usepackage{siunitx}
\usepackage{xcolor}
\usepackage[colorlinks=true,linkcolor=blue,citecolor=blue,urlcolor=blue]{hyperref}

\captionsetup{font=small,labelfont=bf}

% ------------------------------------------------------------------
% Figure helper: shows the image if exported, otherwise a labelled
% placeholder box, so the document always compiles. Drop the PNGs in
% ./figures/ and the boxes are replaced automatically.
% ------------------------------------------------------------------
\newcommand{\fig}[2][\linewidth]{%
  \IfFileExists{figures/#2}%
    {\includegraphics[width=#1]{figures/#2}}%
    {\fbox{\parbox[c][3.4cm][c]{0.95\linewidth}{\centering\ttfamily
       figures/#2\\[2pt]\normalfont\small(export this plot from the R notebook)}}}%
}

% ==================================================================
% KEY NUMERICAL RESULTS -- SINGLE SOURCE OF TRUTH
% These come from the exploratory run (n = 400 subsample, 2 chains,
% 1000 burn-in + 2000 iterations). Re-run the notebook on the full
% n = 2000 with more chains/iterations and update ONLY this block.
% ==================================================================
\newcommand{\Nfull}{2000}
\newcommand{\Nfit}{400}
\newcommand{\nchains}{2}
\newcommand{\nburn}{1000}
\newcommand{\niter}{2000}

% data summary (full n = 2000)
\newcommand{\lfmin}{0.426}
\newcommand{\lfmed}{0.733}
\newcommand{\lfmean}{0.724}
\newcommand{\lfmax}{0.999}

% Task 1 -- chosen model
\newcommand{\bestH}{3}
\newcommand{\waicHthree}{-567.1}
\newcommand{\seHthree}{22.6}

% Task 1 -- components (label-switch corrected, sorted by mean)
\newcommand{\piA}{0.367}\newcommand{\alA}{40.5}\newcommand{\beA}{31.6}\newcommand{\muA}{0.562}
\newcommand{\piB}{0.373}\newcommand{\alB}{66.0}\newcommand{\beB}{22.2}\newcommand{\muB}{0.749}
\newcommand{\piC}{0.261}\newcommand{\alC}{8.16}\newcommand{\beC}{0.64}\newcommand{\muC}{0.927}

% Task 2 -- predictive modes
\newcommand{\modeOne}{0.564}
\newcommand{\modeTwo}{0.756}

% diagnostics
\newcommand{\rhatlik}{1.02}
\newcommand{\esslik}{109}

% Task 4 -- harmonics
\newcommand{\bestK}{3}

% prediction exercise
\newcommand{\ntest}{80}
\newcommand{\logloss}{0.295}
\newcommand{\loglossbase}{1.060}
\newcommand{\acc}{0.900}
\newcommand{\accbase}{0.384}
% ==================================================================

\title{\textbf{Bayesian Analysis of the Hourly City Load Factor}\\[4pt]
\large Beta mixture modelling, latent structure and time-of-day prediction\\
\normalsize Bayesian Learning and Monte Carlo Simulations --- Final Project 2026 (Dataset 11)}
\author{}
\date{\today}

\begin{document}
\maketitle

\begin{abstract}
\noindent
We model the distribution of the \emph{hourly city load factor} of Tetouan
(Morocco) --- the ratio between the average electrical load in a given hour
and the daily peak load --- as a finite mixture of Beta densities estimated
by MCMC in JAGS. Model selection by WAIC over $H=1,\dots,5$ components
selects $H=\bestH$, which recovers the multimodal shape of the empirical
distribution that a single Beta cannot reproduce. We resolve the
label-switching that affects mixture posteriors, report identifiable
per-component summaries with $95\%$ credible intervals, and confirm
robustness through a prior sensitivity analysis. We then link the latent
load regimes to the time of day: a maximum-a-posteriori (MAP) cluster
assignment shows a clear diurnal pattern, and a multinomial logistic
regression on a Fourier expansion of \texttt{hour} reproduces it smoothly.
An out-of-sample exercise confirms that the hour of the day is a genuinely
predictive covariate of the load regime.
\end{abstract}

% ==================================================================
\section{Data and problem description}
% ==================================================================

\textbf{Source and construction.} The data derive from Salam \& El~Hibaoui,
\emph{Power Consumption of Tetouan City} (UCI ML Repository, dataset 849),
which records the active power drawn by the three distribution zones of
Tetouan during 2017, sampled every ten minutes by the utility's SCADA
system. For each (day, hour) cell the city-wide average power
$\overline P(d,h)=\mathrm{Zone1}+\mathrm{Zone2}+\mathrm{Zone3}$ was computed,
and then the \emph{hourly load factor}
\begin{equation}
\mathrm{LF}(d,h)=\frac{\overline P(d,h)}{\max_{h'}\overline P(d,h')}\in(0,1],
\end{equation}
i.e.\ the average load in hour $h$ relative to the highest hourly average of
the same day. By construction the day-peak hour attains $1$ and the deepest
off-peak hour the smallest value. The shared file is a uniform random
subsample of $\Nfull$ (day, hour) cells, with values clamped to
$(10^{-3},\,1-10^{-3})$ so that the Beta log-density is finite at the
boundary. The variables are \texttt{date} (calendar date),
\texttt{hour}~$\in\{0,\dots,23\}$ and \texttt{load\_factor}.

\textbf{Empirical features.} The load factor is a continuous proportion
confined to the unit interval, with summary statistics (full sample)
$\min=\lfmin$, median $=\lfmed$, mean $=\lfmean$, $\max=\lfmax$. The
empirical density (Figure~\ref{fig:hist}) is clearly \emph{not} unimodal:
it shows a heavier off-peak bulk, an intermediate shoulder and a sharp
accumulation of mass near $1$ corresponding to hours close to the daily
peak. This bounded, skewed and multimodal shape is precisely what motivates
a \emph{mixture of Beta distributions} rather than a single Beta or a
Gaussian model.

\textbf{Goals.} Following the project brief we (1) compare a single Beta
against Beta mixtures with $H=1,\dots,5$ components and select $H$ by WAIC;
(2) overlay the posterior predictive density on the histogram and check that
the modes are recovered; (3) assign each observation to its MAP component
and study how component membership varies with the hour of the day; and
(4) fit a multinomial logistic regression that predicts the component label
from a Fourier expansion of \texttt{hour}, comparing it with the empirical
fractions and validating it out of sample.

\begin{figure}[H]
\centering
\fig[0.7\linewidth]{01_hist.png}
\caption{Empirical density of the hourly load factor ($n=\Nfull$). The
bounded, right-skewed and multimodal profile motivates a Beta-mixture
likelihood.}
\label{fig:hist}
\end{figure}

% ==================================================================
\section{Model specification}
% ==================================================================

\subsection{Likelihood: finite Beta mixture}

We treat the load factors $x_1,\dots,x_N$ as conditionally i.i.d.\ draws
from a convex combination of $H$ Beta densities,
\begin{equation}
f(x\mid\bm\pi,\bm\alpha,\bm\beta)=\sum_{h=1}^{H}\pi_h\,
\mathrm{Beta}(x\mid\alpha_h,\beta_h),
\qquad
\mathrm{Beta}(x\mid\alpha,\beta)=\frac{x^{\alpha-1}(1-x)^{\beta-1}}{B(\alpha,\beta)},
\label{eq:mix}
\end{equation}
with mixing weights $\pi_h\ge 0$, $\sum_h\pi_h=1$ and strictly positive
shape parameters $\alpha_h,\beta_h>0$. The Beta family is the natural choice
on $(0,1)$: each component has mean $\mu_h=\alpha_h/(\alpha_h+\beta_h)$ and
can be made left-skewed, right-skewed, symmetric or sharply peaked depending
on $(\alpha_h,\beta_h)$, so a small number of components can capture distinct
``load regimes'' (off-peak, intermediate, near-peak).

\subsection{Latent-variable (data-augmentation) formulation}

For inference in JAGS we augment the model with a latent allocation variable
$z_i\in\{1,\dots,H\}$ for each observation, recovering \eqref{eq:mix} on
marginalisation:
\begin{equation}
z_i\mid\bm\pi\;\sim\;\mathrm{Categorical}(\bm\pi),
\qquad
x_i\mid z_i\;\sim\;\mathrm{Beta}(\alpha_{z_i},\beta_{z_i}).
\end{equation}
This is the standard mixture data augmentation: conditioning on $z_i$ turns
the likelihood into a product of ordinary Beta terms that JAGS samples
efficiently with a Gibbs/slice scheme.

\subsection{Prior distributions and hyperparameters}

We use weakly informative, conditionally conjugate priors that let the data
drive the posterior:
\begin{align}
\bm\pi &\sim \mathrm{Dirichlet}(\mathbf 1_H), &&\text{(symmetric, flat on the simplex)}\\
\alpha_h &\sim \mathrm{Exponential}(\lambda), &
\beta_h &\sim \mathrm{Exponential}(\lambda), && h=1,\dots,H,
\end{align}
with rate $\lambda=0.1$. The symmetric Dirichlet$(\mathbf 1_H)$ is uniform
over the probability simplex and expresses no preference among components.
The Exponential prior keeps the shape parameters strictly positive while
remaining diffuse: with $\lambda=0.1$ the prior mean of each shape parameter
is $1/\lambda=10$ and the prior standard deviation is also $10$, comfortably
covering the range of shapes the data require (posterior $\alpha,\beta$ up to
$\sim\!70$). The influence of these two hyperparameters ($\lambda$ and the
Dirichlet concentration $c$) is assessed explicitly in
Section~\ref{sec:sens}.

\subsection{Posterior computation and model-selection criterion}
\label{sec:algo}

\textbf{Algorithm.} All posteriors are sampled by MCMC in \texttt{JAGS}
(via \texttt{rjags}). For the results reported here we ran $\nchains$ chains
with dispersed initial values, $\nburn$ burn-in iterations and $\niter$
retained iterations per chain. Convergence is monitored with trace plots,
the Gelman--Rubin statistic $\widehat R$ and effective sample size
(Section~\ref{sec:diag} and Appendix~\ref{app:diag}). The present figures
are produced on a $\Nfit$-row exploratory subsample for speed; the final
report will be regenerated on the full $n=\Nfull$ sample with additional
chains and iterations.

\textbf{Model selection.} Competing values of $H$ are compared with the
\emph{Widely Applicable Information Criterion} (WAIC), a fully Bayesian
estimate of out-of-sample predictive accuracy,
\begin{equation}
\mathrm{WAIC}=-2\bigl(\mathrm{lppd}-p_{\mathrm{WAIC}}\bigr),\quad
\mathrm{lppd}=\sum_{i=1}^N\log\Bigl(\tfrac1S\sum_{s=1}^S f(x_i\mid\bm\theta^{(s)})\Bigr),\quad
p_{\mathrm{WAIC}}=\sum_{i=1}^N\mathrm{Var}_{s}\!\bigl[\log f(x_i\mid\bm\theta^{(s)})\bigr],
\end{equation}
where $\bm\theta^{(s)}$ is the $s$-th posterior draw and $p_{\mathrm{WAIC}}$
penalises effective complexity. \textbf{Crucially, for a mixture WAIC must
use the \emph{marginal} per-observation density}
$\log\sum_h \pi_h^{(s)}\mathrm{Beta}(x_i\mid\alpha_h^{(s)},\beta_h^{(s)})$,
integrating out $z_i$. We compute this quantity inside JAGS with a numerically
stable log-sum-exp evaluation and feed it to \texttt{loo::waic}. Using the
density conditional on the sampled $z_i$ would answer a different predictive
question, inflate $p_{\mathrm{WAIC}}$ and spuriously favour more components;
the marginal density is also invariant to label permutations, so this WAIC
is immune to label switching.

% ==================================================================
\section{Model selection: how many components?}
% ==================================================================

We fit the single Beta ($H=1$) and the mixtures $H=2,\dots,5$ and rank them
by WAIC, reading the criterion together with its standard error (SE) rather
than picking the bare minimum: two models whose WAIC differ by less than
$\sim\!1$--$2$ SE are statistically indistinguishable, and among tied models
the most parsimonious is preferable.

\begin{table}[H]
\centering
\caption{Model selection by WAIC. \texttt{ELPD} is the expected log
predictive density, \texttt{p\_eff}$=p_{\mathrm{WAIC}}$ the effective number
of parameters, lower WAIC is better. The right block reports
\texttt{loo\_compare} differences against the best model.}
\label{tab:waic}
\begin{tabular}{lrrrr@{\hskip 2.2em}rr}
\toprule
Model & ELPD & p\_eff & WAIC & SE & $\Delta$elpd & se\_diff\\
\midrule
$H=1$ & 163.9 & 2.7 & $-327.8$ & 14.4 & $-119.7$ & 12.8\\
$H=2$ & 249.7 & 3.9 & $-499.4$ & 20.5 & $-33.8$  & 6.7\\
$H=\mathbf{3}$ & \textbf{283.5} & 6.2 & $\mathbf{-567.1}$ & 22.6 & \textbf{0.0} & 0.0\\
$H=4$ & 282.3 & 6.6 & $-564.6$ & 22.1 & $-1.2$ & 0.5\\
$H=5$ & 283.3 & 6.4 & $-566.7$ & 22.4 & $-0.2$ & 0.4\\
\bottomrule
\end{tabular}
\end{table}

WAIC drops sharply from $H=1$ to $H=3$ ($\Delta\mathrm{elpd}\approx 120$ and
$34$ standard-error-significant units for the two steps) and then plateaus:
$H=4$ and $H=5$ are within one SE of $H=3$ and do not improve the fit
($\Delta\mathrm{elpd}=-1.2$ and $-0.2$, both well inside their se\_diff).
The lowest-WAIC \emph{and} the most parsimonious choice therefore coincide
at $H=\bestH$ (Figure~\ref{fig:waicH}). We adopt $H=\bestH$ for all
subsequent analysis. The effective number of parameters
($p_{\mathrm{WAIC}}\approx6.2$ for three components) is close to the
$3H-1=8$ free parameters, indicating that all three components are genuinely
supported by the data rather than redundant.

\begin{figure}[H]
\centering
\fig[0.7\linewidth]{02_waic_H.png}
\caption{WAIC against the number of components $H$ with $\pm1$~SE bars. The
curve falls steeply until $H=\bestH$ and then flattens; the simplest model
within one SE of the minimum is $H=\bestH$.}
\label{fig:waicH}
\end{figure}

% ==================================================================
\section{Posterior analysis and interpretation}
% ==================================================================

\subsection{Label switching and identifiability}
\label{sec:diag}

In a mixture the components are exchangeable: permuting their labels leaves
the likelihood unchanged, so the posterior has up to $H!$ equivalent modes.
When independent chains explore different permutations the raw parameter
traces cross and swap --- \emph{label switching} --- and naively averaging
component parameters across draws blends different permutations. The trace
plots and the per-component $\widehat R$ values (several well above the
conventional $1.10$ threshold; Appendix~\ref{app:diag}) display exactly this
signature.

This is, however, a \emph{labelling} artefact and not a failure of
convergence. Evaluating $\widehat R$ and ESS on the
\emph{permutation-invariant} marginal log-likelihood
$\ell^{(s)}=\sum_i\log f(x_i\mid\bm\theta^{(s)})$ gives
$\widehat R=\rhatlik$ and a healthy ESS$\,\approx\esslik$, confirming that
the sampler has reached the global posterior. To report identifiable
per-component summaries we apply a \textbf{post-hoc relabelling}: at every
draw the $H$ components are sorted by their Beta mean
$\mu_h=\alpha_h/(\alpha_h+\beta_h)$, so ``Component~1'' is consistently the
lowest-mean (off-peak) regime and ``Component~3'' the highest-mean
(near-peak) regime. This does not alter the fit --- the mixture density and
WAIC are unchanged --- it only fixes the labels for interpretation.

\subsection{Posterior summaries of the selected model}

Table~\ref{tab:comp} reports, for the $H=\bestH$ model, the posterior mean
weight, shape parameters, implied component mean and $95\%$ equal-tailed
credible intervals after relabelling.

\begin{table}[H]
\centering
\caption{Posterior summaries of the $H=\bestH$ Beta mixture (posterior
means; $95\%$ credible intervals for the shape parameters). Components are
ordered by increasing mean.}
\label{tab:comp}
\begin{tabular}{clccccc}
\toprule
Comp. & Regime & $\pi_h$ & $\alpha_h$ & $\beta_h$ & $\mu_h$ &
$\alpha_h$ CI\,/\,$\beta_h$ CI\\
\midrule
1 & off-peak / low   & \piA & \alA & \beA & \muA & $[29.4,\,52.1]$ / $[22.3,\,41.6]$\\
2 & intermediate     & \piB & \alB & \beB & \muB & $[43.1,\,111.7]$ / $[14.8,\,36.9]$\\
3 & near-peak / high & \piC & \alC & \beC & \muC & $[4.3,\,14.3]$ / $[0.5,\,0.8]$\\
\bottomrule
\end{tabular}
\end{table}

The three components admit a clean physical reading. Component~1
($\mu\approx\muA$, weight $\approx\piA$) is the broad off-peak regime;
Component~2 ($\mu\approx\muB$, weight $\approx\piB$) is the intermediate
load regime and is the most concentrated (largest $\alpha$); Component~3
($\mu\approx\muC$, weight $\approx\piC$) captures the sharp accumulation of
mass near the daily peak --- its very small $\beta\approx\beC$ produces the
strong left build-up towards $1$. The credible intervals are away from zero
for all shape parameters, confirming that each regime is well identified.

\subsection{Posterior predictive check and mode recovery}

The fully Bayesian overlay is the \emph{posterior predictive density},
obtained by averaging the mixture density over the entire posterior,
\begin{equation}
p(\tilde x\mid\text{data})=\frac1S\sum_{s=1}^S\sum_{h=1}^H
\pi_h^{(s)}\,\mathrm{Beta}\!\bigl(\tilde x\mid\alpha_h^{(s)},\beta_h^{(s)}\bigr).
\end{equation}
We average \emph{densities}, not parameters, so the curve propagates
posterior uncertainty and is automatically invariant to label switching.
Figure~\ref{fig:anatomy} overlays this density and its weighted components
on the histogram; Figure~\ref{fig:ppd} marks its local maxima. The
predictive modes fall at $\modeOne$ and $\modeTwo$, aligning with the
visible peaks of the empirical distribution (the near-peak component
contributes a shoulder rather than a separated mode because its mass piles
against the upper boundary). The $H=\bestH$ mixture thus reproduces a
multimodal profile that a single Beta is structurally unable to capture.

\begin{figure}[H]
\centering
\begin{subfigure}{0.49\linewidth}\centering
\fig{03_anatomy.png}
\caption{Mixture anatomy: overall predictive density (solid) and weighted
components (dashed).}
\label{fig:anatomy}
\end{subfigure}\hfill
\begin{subfigure}{0.49\linewidth}\centering
\fig{04_ppd_modes.png}
\caption{Posterior predictive overlay with recovered modes at $\modeOne$ and
$\modeTwo$.}
\label{fig:ppd}
\end{subfigure}
\caption{Posterior predictive density of the selected $H=\bestH$ model
against the empirical histogram.}
\end{figure}

\subsection{Prior sensitivity analysis}
\label{sec:sens}

The conclusions rest on two hyperparameters: the Exponential rate $\lambda$
on the shape parameters and the Dirichlet concentration $c$ on the weights.
We refit the $H=\bestH$ model under four settings, varying one at a time
(Table~\ref{tab:sens}).

\begin{table}[H]
\centering
\caption{Prior sensitivity for the $H=\bestH$ model. WAIC and effective
complexity under four hyperparameter settings.}
\label{tab:sens}
\begin{tabular}{lccrrr}
\toprule
Setting & $\lambda$ & $c$ & ELPD & p\_eff & WAIC\\
\midrule
Reference         & 0.10 & 1.0 & 283.5 & 6.2  & $-567.1$\\
Very diffuse      & 0.01 & 1.0 & 284.5 & 7.3  & $-569.0$\\
Informative       & 1.00 & 1.0 & 247.3 & 3.0  & $-494.7$\\
Sparse Dirichlet  & 0.10 & 0.1 & 274.1 & 12.5 & $-548.2$\\
\bottomrule
\end{tabular}
\end{table}

The reference and the very diffuse priors give essentially identical WAIC
and overlapping posterior predictive curves (Figure~\ref{fig:sens}),
confirming that under weak priors the inference is data-driven. Two settings
are informative enough to matter. The tight Exponential ($\lambda=1$, prior
mean shape $=1$) pulls the shapes towards $1$, cannot represent the
concentrated intermediate and near-peak regimes, and noticeably degrades
predictive accuracy (WAIC $-494.7$, $p_{\mathrm{WAIC}}$ collapsing to $3$).
The sparse Dirichlet ($c=0.1$) encourages a single dominant component and
also worsens the fit. The analysis therefore both validates the reference
prior and shows where an over-confident prior would distort the result ---
exactly the diagnostic value a sensitivity study should provide.

\begin{figure}[H]
\centering
\fig[0.7\linewidth]{05_sensitivity.png}
\caption{Posterior predictive density of the $H=\bestH$ model under the four
prior settings of Table~\ref{tab:sens}. The reference and very-diffuse
curves coincide; the informative-shape and sparse-Dirichlet priors depart.}
\label{fig:sens}
\end{figure}

% ==================================================================
\section{Latent load regimes across the day (post-hoc exploration)}
% ==================================================================

Having identified three load regimes from the load factor alone, we ask how
they relate to the time of day. Each observation is assigned to its MAP
component using posterior \emph{responsibilities} computed on the
mean-ordered components,
\begin{equation}
r_{ih}=\frac1S\sum_{s=1}^S
\frac{\pi_{(h)}^{(s)}\mathrm{Beta}(x_i\mid\alpha_{(h)}^{(s)},\beta_{(h)}^{(s)})}
{\sum_k\pi_{(k)}^{(s)}\mathrm{Beta}(x_i\mid\alpha_{(k)}^{(s)},\beta_{(k)}^{(s)})},
\qquad \hat z_i=\arg\max_h r_{ih},
\end{equation}
where $(h)$ denotes the $h$-th component by increasing mean, which keeps the
assignment stable across chains. Plotting the empirical fraction of
observations in each component against \texttt{hour}
(Figure~\ref{fig:byhour}) reveals a strong, smooth diurnal cycle: the
low-mean (off-peak) component dominates at night, the intermediate component
prevails through the working day, and the high-mean (near-peak) component
concentrates around the daily peak hours. The latent structure recovered
\emph{unsupervised} from the load factor is thus tightly organised by the
hour of the day --- a substantive validation that the components correspond
to real operating regimes rather than statistical artefacts.

\begin{figure}[H]
\centering
\fig[0.7\linewidth]{06_byhour.png}
\caption{Empirical fraction of observations in each MAP component as a
function of \texttt{hour}. Off-peak (Comp.~1) dominates at night,
intermediate (Comp.~2) during the day, near-peak (Comp.~3) around the daily
peak.}
\label{fig:byhour}
\end{figure}

% ==================================================================
\section{Predicting the load regime from the hour: multinomial logit}
% ==================================================================

We now turn the relationship around and ask \emph{how predictable the
regime is from the time of day}, treating the MAP label $\hat z_i$ as the
response of a supervised model whose only covariate is \texttt{hour}.
Because \texttt{hour} is cyclic ($23$ is adjacent to $0$), we expand it on a
Fourier basis of period $24$,
\begin{equation}
s_{ik}=\sin\!\Bigl(\tfrac{2\pi k\,\mathrm{hour}_i}{24}\Bigr),\qquad
c_{ik}=\cos\!\Bigl(\tfrac{2\pi k\,\mathrm{hour}_i}{24}\Bigr),\qquad k=1,\dots,K,
\end{equation}
and use these $2K$ functions in a multinomial logistic regression with the
lowest-mean component as reference,
\begin{equation}
\eta_{i,h}=\beta_{0,h}+\sum_{k=1}^K\bigl(\alpha_{k,h}s_{ik}+\gamma_{k,h}c_{ik}\bigr),
\qquad
P(z_i=h)=\frac{e^{\eta_{i,h}}}{\sum_{j=1}^H e^{\eta_{i,j}}},
\end{equation}
with $\beta_{0,1}=\alpha_{k,1}=\gamma_{k,1}=0$. The softmax guarantees
$\sum_h P(z_i=h)=1$ and the predictions are periodic in \texttt{hour} by
construction. We place vague $\mathcal N(0,\,\mathrm{var}=100)$ priors on the
$H-1$ free coefficient blocks and fit in JAGS with the same sampler
settings. The number of harmonics $K$ is selected by WAIC exactly as $H$
was: WAIC keeps improving up to $K=5$ but $K=\bestK$ is already within one SE
of the minimum, so we adopt the parsimonious $K=\bestK$
(Appendix~\ref{app:diag}, Table~\ref{tab:waicK}).

The predicted probability curves (Figure~\ref{fig:task4}) closely track the
empirical fractions of the previous section: both estimate the same quantity
$P(\text{component}\mid\text{hour})$, but the logit borrows strength across
adjacent hours through the Fourier basis and yields a smooth, strictly
periodic profile, whereas the empirical fractions are noisier at hours with
few sampled observations. The agreement confirms that the latent load
regimes are, to a good approximation, a deterministic periodic function of
the time of day.

\begin{figure}[H]
\centering
\fig[0.75\linewidth]{07_task4.png}
\caption{Multinomial logit ($K=\bestK$): predicted $P(\text{component}\mid
\text{hour})$ (curves) overlaid on the empirical fractions of
Figure~\ref{fig:byhour} (points). The smooth periodic model reproduces the
empirical pattern.}
\label{fig:task4}
\end{figure}

% ==================================================================
\section{Prediction exercise (out-of-sample)}
% ==================================================================

The WAIC above is an in-sample criterion. To verify that \texttt{hour}
carries genuine predictive information we split the data into an $80\%$
training and a $20\%$ test set ($n_\text{test}=\ntest$), refit the
$K=\bestK$ logit on the training set only, and predict the regime of each
test observation from its hour. We report the test \emph{log-loss}
$-\frac1{n_\text{test}}\sum_{i}\log\hat P(z_i=\hat z_i\mid\mathrm{hour}_i)$
and the classification accuracy, against a null baseline that ignores
\texttt{hour} and always predicts the training class proportions
(Table~\ref{tab:pred}).

\begin{table}[H]
\centering
\caption{Out-of-sample performance on the held-out test set
($n_\text{test}=\ntest$), against an hour-agnostic baseline.}
\label{tab:pred}
\begin{tabular}{lcc}
\toprule
Metric & Logit model & Baseline\\
\midrule
Log-loss (lower better) & \logloss & \loglossbase\\
Accuracy (higher better) & \acc & \accbase\\
\bottomrule
\end{tabular}
\end{table}

The model more than halves the log-loss ($\logloss$ vs $\loglossbase$) and
raises accuracy from $\accbase$ to $\acc$. Knowing the time of day therefore
substantially improves the prediction of the load regime over the
unconditional frequencies, and the periodic pattern learned on $80\%$ of the
data reproduces the full-data fractions --- evidence that it is a structural
feature of city load, not an artefact of the particular training split.

% ==================================================================
\section{Conclusions}
% ==================================================================

A three-component Beta mixture provides a parsimonious and interpretable
description of the hourly city load factor of Tetouan. WAIC clearly prefers
the mixture over a single Beta and selects $H=\bestH$ components, which
recover the multimodal empirical density and admit a direct reading as
off-peak, intermediate and near-peak load regimes. We addressed the
label-switching intrinsic to mixtures by checking convergence on a
permutation-invariant scale ($\widehat R=\rhatlik$) and relabelling
components by their mean, allowing identifiable per-component summaries with
credible intervals. A prior sensitivity analysis confirmed that the
conclusions are data-driven under weak priors. Linking the latent regimes to
the hour of the day --- first by MAP exploration, then by a Fourier-based
multinomial logit selected by WAIC --- showed a strong, smooth diurnal
cycle, and an out-of-sample exercise confirmed that \texttt{hour} is a
genuinely predictive covariate of the load regime.

\textbf{Limitations and extensions.} The load factor is constructed by
dividing each day by its own peak, which removes day-level scale; a
hierarchical extension could model day- and season-level variation jointly.
The number of components could be estimated rather than selected, e.g.\ with
a Dirichlet-process mixture. Finally, the figures here use a $\Nfit$-row
exploratory run; the final version will be regenerated on the full
$n=\Nfull$ sample with more chains and iterations, which we expect to
sharpen the credible intervals without changing the qualitative conclusions.

% ==================================================================
\appendix
\section{MCMC diagnostics and supplementary results}
\label{app:diag}
% ==================================================================

\textbf{Convergence.} Figure~\ref{fig:trace} shows trace plots for the
component parameters of the $H=\bestH$ model. The chains run on separate
tracks that occasionally cross --- the characteristic label-switching
pattern --- and the corresponding raw per-component $\widehat R$ values
exceed $1.10$ (Table~\ref{tab:rhat}). Evaluated on the permutation-invariant
marginal log-likelihood, however, $\widehat R=\rhatlik$ and
ESS$\,\approx\esslik$, confirming global convergence;
the high raw $\widehat R$ is a labelling artefact, not a mixing failure.

\begin{table}[H]
\centering
\caption{Raw (label-affected) Gelman--Rubin $\widehat R$ for the $H=\bestH$
component parameters. Values above $1.10$ reflect label switching, not
non-convergence (cf.\ the invariant-scale $\widehat R=\rhatlik$).}
\label{tab:rhat}
\begin{tabular}{lc@{\hskip 2em}lc@{\hskip 2em}lc}
\toprule
Par. & $\widehat R$ & Par. & $\widehat R$ & Par. & $\widehat R$\\
\midrule
$\alpha_1$ & 7.48 & $\beta_1$ & 9.46 & $\pi_1$ & 4.86\\
$\alpha_2$ & 1.59 & $\beta_2$ & 1.60 & $\pi_2$ & 1.23\\
$\alpha_3$ & 10.04 & $\beta_3$ & 12.01 & $\pi_3$ & 3.91\\
\bottomrule
\end{tabular}
\end{table}

\begin{figure}[H]
\centering
\fig[0.85\linewidth]{A1_traces.png}
\caption{Trace plots (and marginal densities) for the $H=\bestH$ component
parameters, showing the label-switching signature.}
\label{fig:trace}
\end{figure}

\textbf{Selection of the number of harmonics.} Table~\ref{tab:waicK} reports
WAIC for the multinomial logit over $K=1,\dots,5$ harmonics. WAIC improves
through $K=5$, but $K=\bestK$ is already within one SE of the minimum
($\Delta\mathrm{elpd}=-2.3$, se\_diff $4.5$), so the parsimonious choice is
$K=\bestK$.

\begin{table}[H]
\centering
\caption{WAIC for the Task-4 multinomial logit by number of harmonics $K$.}
\label{tab:waicK}
\begin{tabular}{lrrrr@{\hskip 2.2em}rr}
\toprule
Model & ELPD & p\_eff & WAIC & SE & $\Delta$elpd & se\_diff\\
\midrule
$K=1$ & $-142.6$ & 8.7 & 285.2 & 33.3 & $-30.4$ & 10.2\\
$K=2$ & $-131.9$ & 13.7 & 263.8 & 38.6 & $-19.7$ & 11.3\\
$K=\mathbf{3}$ & $-114.5$ & 13.0 & 229.0 & 28.0 & $-2.3$ & 4.5\\
$K=4$ & $-113.7$ & 15.2 & 227.4 & 28.7 & $-1.5$ & 3.4\\
$K=5$ & $-112.2$ & 16.8 & 224.4 & 27.9 & 0.0 & 0.0\\
\bottomrule
\end{tabular}
\end{table}

\end{document}
>>>>>>> Stashed changes
